\section{Introduction}
\label{introduction}
The Lonely Runner conjecture was first stated by J. M. Wills in \cite{JMWills}, where it was presented as a purely formal problem in the field of Diophantine approximation. It remained so until 1998, where it was given its more memorable name, and intuitive interpretation, from \cite{Bienia97flows.view-obstructions}. As of 2010 the Lonely Runner conjecture has only been proved for up to $n$ = 6 \cite{serra_thelonely}.

There exist several equivalent interpretations of the Lonely Runner conjecture, and I will here state widely used versions\footnote{See \cite{ANote}}:
\begin{enumerate}
\item Imagine $n$ runners at a circular, unit length track\footnote{i.e. the track has a circumference of 1}, where every runner runs with constant, pairwise different speeds\footnote{These speeds can, without loss of generality be assumed to be in $\N$ - see \cite{Bienia97flows.view-obstructions} and Section \ref{integerSpeeds}}, there will then exist a time t where all runners will be at least $\frac{1}{n + 1}$ units away from their common start-point.\\

\item Imagine we have $n$ + 1 runners, then the conjecture states that for every runner there exists a time where all other runners are at least $\frac{1}{n + 1}$ units away from it. Everything else is the same as in the first interpretation \cite{Bienia97flows.view-obstructions}.\\
\end{enumerate}

More formally the problem can be stated thus: 
Given any $n$ positive integers $w_1, w_2, \ldots, w_n$, there exists a real number $x$ such that 
\eqn{
\label{eqa:lonelyRunner} \Vert w_i x\Vert \geq \frac{1}{n+1}
}

for $i = 1, 2, \ldots, n$, where for $y \in \R$, $\Vert y \Vert$ is the distance from $y$ to the nearest integer. I owe the above formulations of the conjecture to \cite{ANote}. In this report it is useful to think of the track as the one-dimensional Torus $\mathbb{T} = [0,1)$.

\subsection{Expectations of the Reader}
I expect the reader to have the mathematical knowledge required to understand both the formal and informal description of the Lonely Runner Conjecture, as well having an understanding of university level Computational Geometry\footnote{For instance having read \cite{citeulike:3347056} or similar}.

\subsection{Scope and Limitations}
\begin{itemize}
\item In this project I will not attempt to prove the Lonely Runner conjecture for all n. 
\item I will create a program that can verify whether a given configuration of runners is true for equation \eqaref{eqa:lonelyRunner}, and make it return a result, within a reasonable time, for up to 1000 runners. The final program will work with more than 1000 runners, but I will prioritise readability, maintainability and portability of the program instead of extreme optimisations. 
\item The program will include a system to test a range of runner configurations against equation \eqaref{eqa:lonelyRunner}.
\item I will only focus on runners with integer speeds - see \ref{integerSpeeds} for an justification of this.
\end{itemize}

\subsection{Computer Specifications}
\label{specs}
I would like to define which kind of computer I would require for running my program. As mentioned above, the Lonely Runner conjecture has direct applications for View-Obstruction problems and colouring of regular chromatic graphs. Researchers working with in these fields have no special reason to work with exceedingly powerful computers, or computers with special hardware.
 
Based on the above I would argue that any algorithm I devise and implement should be able to run on a standard computer with no special hardware\footnote{Say hardware dedicated to fast floating point calculations}. I will therefore design my verification program to run on a 2 Ghz computer with the recommended amount of RAM for the computer's OS. 

I reserve the right to use third-party software libraries - as long as these can be freely distributed and installed.

\subsection{Integer Speeds}
\label{integerSpeeds}
In this section I will prove that my program will only have to work with integer speeds.
Let us assume for a given configuration of runner speeds T, there is a non-empty set $S \subset T$, of runner speeds that are not integers. Since we are dealing with the speeds of runners, it is clear that there are 3 cases we have to worry about: 1) Either all elements in S belong to $\mathbb{Q}$, 2) All elements are irrational numbers, 3) There is a mix of irrational and rational numbers.

I will now prove that in all 3 cases we only have to worry about finding a solution for the speeds that can be reduced to integers.

\begin{description}
\item[Only rational numbers in S:] In this instance we can convert all the runner speeds into integers, by multiplying all speeds in T with the product of the denominators of the rational numbers in S.
\item[Only irrational numbers in S, and S = T:]
Since we are only dealing with irrational numbers, then from example 6.1 in \cite{uniform} they must be uniform distribution mod 1 in $\R$. 
From Definition 6.1 in \cite{uniform}, we furthermore have that for every pair $\mathbf{a}, \mathbf{b} \in \R^{|S|}$, where $\mathbf{a} = (a, a, \ldots, a)$, and $\mathbf{b} = (b, b, \ldots, b)$ with $\mathbf{0} \leq \mathbf{a} < \mathbf{b} \leq \mathbf{1}$, and every $s = (s_1, s_2, \ldots, s_n)$, where $sk = (s_1k, s_2k, \ldots, s_nk)$ we have

$$
\lim_{N \rightarrow \infty} \frac{
|\E{
sn | \E{sn} \in [\mathbf{a}, \mathbf{b}) \;\mathrm{for}\;n \leq N
}|
}{
N
} = 
\prod^{|S|}_{j=1}(b_j - a_j) 
$$

where $\E{x} = x - [x]$, where $[x]$ is the largest integer $y$ less than or equal to x, and for $[s] = ([s_1], [s_2], \ldots, [s_n])$.

From this it is clear that for any interval $[a, b]$, it will be possible to find a $k \in \N$ large enough to place all the runners inside it. Thus we pick a $k$ such that all runners in T are inside an interval that makes equation \eqaref{eqa:lonelyRunner} hold. So it holds trivially if we only have irrational numbers.

\item[T is a mixture of natural/rational numbers and irrational numbers:]
In this case we first multiply all the speeds in T with the product of all the denominator of the fractions in T, giving us the set of speeds $T\prime$. Let $N\prime$ be the true subset of $T\prime$ that only contains integer speeds, and $R\prime$ be the true subset of $T\prime$ that only contains irrational speeds, so we have $T\prime = N\prime \cup R\prime$, and $N\prime \cap R\prime = \emptyset$.

Let us assume that in that case there must exist a real number $t$, which makes equation \eqaref{eqa:lonelyRunner} hold for $N\prime$ (where the number of runners is $|N\prime|$). Let us observe the vector $(r_1, r_2, \ldots, r_k, n_1, n_2, \ldots, n_m)$, where $r_i \in R\prime$ and $|R\prime| = k$, and $n_i \in N\prime$ and $|N\prime| = m$. 

From the last case it is clear that there must exist an integer which makes equation \eqaref{eqa:lonelyRunner} hold for all the speeds in $R\prime$, so observing $t$ from before, let us now choose a $q \in \N$, such that for all the speeds in $R\prime$, $x = q + t$ makes them hold for equation \eqaref{eqa:lonelyRunner}\footnote{Since we are able to cram all our runners into an arbitrary small interval anywhere on $\mathbb{T}$, this can be done}.

So now if observe the vector 
\begin{equation}
\begin{split}
&(q + t)(r_1, r_2, \ldots, r_k, n_1, n_2, \ldots, n_m) =\\ 
&((q + t)r_1, (q + t)r_2, \ldots, (q + t)r_n, (q + t)n_1, (q + t)n_2, \ldots, (q + t)n_m) 
\end{split}
\end{equation}

It might not be not entirely clear that that the natural numbers in the above vector gives us a solution, so let us take a closer look: 
$(q + t)n_i = q * n_i + t * n_i$, and $\Vert q * n_i + t * n_i \Vert = \Vert t * n_i \Vert$ since $q \in \N$, and since we started by assuming $t$ made all speeds in $N\prime$ hold with regards to equation \eqaref{eqa:lonelyRunner}, it follow by tautology that $\Vert t * n_1 \Vert \geq \frac{1}{n+1}$ for all $n_i \in N\prime$.

Therefore if there exists a solution for integer speeds, there must also exist a solution for the irrational numbers, and we thus we only have to worry about finding a solution for the integer speeds.
\end{description}

\subsection{Terminology}
\label{Termonolgy}
\begin{figure}[H]
  \centering
  \includegraphics[width=0.3\textwidth, angle=90]{./images/circleZoneEPS}
  \caption{\label{circleZoneImg}An illustration of the Runner track. The greyed out part is the Zone}
\end{figure}

In this report I will refer to runners as male for simplicity.

\begin{defi}
\label{def:theZone}
In this report I will refer to the track interval [$\frac{1}{n + 1}$, $\frac{n}{n+1}$] as the Zone, and a runner who is in this interval, as ``being in the Zone''. see figure \ref{circleZoneImg} for an illustration of the Zone.
\end{defi}

\begin{defi}
\label{def:config}
Configurations of runners, will in this report refer to a set S of $n \in \N$ runners, where for all runners $r, r\prime \in S$, $r \neq r\prime$, $r_{speed}, r\prime_{speed} \in \N$, and $r_{speed} \neq r\prime_{speed}$.
\end{defi}

\hide{
\begin{defi}
\label{def:fakeRunner}
When talking about the fake runner, I will be referring to a runner that is not actually in race, but rather a runner with speed 1, who I have introduced in order to have a convenient termination criterion for the Geometrical algorithm\footnote{\pref{compGeo}}. \pref{termination} for more.
\end{defi}  
}
\subsection{Canonical Example}
Through the rest of this report I will use the following configuration
of runners to give an example:
We have 4 runners with the speeds of 3, 5, 6, 9

\subsection{Background Material Used}
\label{background}
For this project I have read the following papers dealing with the Lonely Runner conjecture: ``View-Obstruction Problems''\cite{Bienia97flows.view-obstructions}, ``The lonely runner problem with seven runners'' \cite{serra_thelonely}, ``Regular chromatic number and the lonely runner problem'' \cite{Barajas2007479}, ``View-Obstruction Problems'' \cite{springerlink:10.1007/BF01832623}, ``Tight Instances of the Lonely Runner'' \cite{Goddyn96tightinstances}, ``A Note on the Lonely Runner Conjecture'' \cite{ANote} and ``Invisible runners in finite fields'' \cite{invis}.

I have also used the follow support literature:
``Computational Geometry'' \cite{citeulike:3347056}, ``Uniform Distribution of Sequences'' \cite{uniform} and ``Kalkulus'' \cite{kalkulus}.

\subsection{Source Code and Test Data}
The source code for the program can be found in my GitHub repo at \underline{https://github.com/kjmikkel/LonelyRunner}, there you will also be able to find the machine readable data I have used as the basis for my results, as well as the source code for this report.

\subsection{Overview}
In this section I give a quick overview of what the different sections in report will cover.
\begin{description}
\item[Section \ref{introduction}] The introduction to the report, which will introduce the subject and explain my goals
\item[Section \ref{choiceOfMethod}] A discussion of the two different approaches that can decide whether or not the Lonely Runner conjecture holds for a given configuration
\item[Section \ref{implementation}] Description of the implementation details and an analysis of the scalability of the implementation
\item[Section \ref{test}] Description of the tests I am going to perform.
\item[Section \ref{results}] Interpretation of the results gained by running the tests described in the previous section.
\item[Section \ref{conclusion}] The conclusion of the report, where I will sum up my findings and discuss future work.
\end{description}
